\mode<presentation>
{
 \usetheme{JuanLesPins}
 \usefonttheme{serif}
 \usecolortheme{beaver}
 \setbeamercovered{invisible} \setbeamertemplate{blocks}[rounded][shadow=true] 
 \setbeamertemplate{navigation symbols}{} 
 \setbeamertemplate{footline}[frame number]
 \usecolortheme[RGB={122,4,24}]{structure}
}

\setcounter{tocdepth}{1}

%%%%%%%%%%%%%%%%%%%%%%%%%%%%%%%%%%%%%%%%%%%
%%%%%%%%%%%%%%%%%%%%%%%%%%%%%%%%%%%%%%%%%%%
%%%%%%%%%%%%%%%%% PACKAGES %%%%%%%%%%%%%%%%%

% BASICAO

\usepackage{lmodern}
\usepackage[T1]{fontenc}
\usepackage[useregional]{datetime2}

% Para usar com PDFTEX
%\usepackage[brazilian]{babel}
%\usepackage[utf8]{inputenc}

% Para usar com XELATEX
\usepackage{polyglossia}
\setdefaultlanguage{brazil}


% FIGURAS
\usepackage{graphicx,float}
\usepackage{booktabs, longtable}


% TABELAS E DIAGRAMAÇÃO
\usepackage{enumerate}
\usepackage{multirow,multicol}
\usepackage{makecell}
\usepackage{array}
%\usepackage{fancybox}
%\usepackage{colortbl}


% ?????????????

% Fonte Courier?
%\usepackage{courier}

% Simbolos
%\usepackage{textcomp}

% Texto caligrafico
%\usepackage{calligra}

% Conexão entre pdfs e dvis
%\usepackage{srcltx}

% Para delimitar environments de tipo float
%\usepackage{cprotect}

% Para texto literal (programação)
%\usepackage{verbatim}

% Pegar números relacionados a referências para uso
%\usepackage{refcount}

% Para colocar relógios de ponteiro
%\usepackage{tdclock}

%%%%%%%%%%%%%%%%%%%%%%%%%%%%%%%%%%%%%%%%%%%
%%%%%%%%%%%%%%%%%%%%%%%%%%%%%%%%%%%%%%%%%%%
%%%%%%%%%%%%%%%%% COISAS MATEMATICAS %%%%%%%%%%%%%%%%%

\usepackage{amsmath,amssymb,amsfonts,xfrac,cancel,bm,mathtools}

% Notação de ângulos para complexos
%\usepackage{steinmetz}

% Notação de derivadas
%\usepackage[thinc]{esdiff}
%\usepackage{commath}

% Setinhas para notação de vetor
%\usepackage{esvect}

% Fonte para notação em rsfs
%\usepackage{mathrsfs}

% Para mapas de karnaugh
%\usepackage{karnaugh-map}


% setup do (circui)tikz
\usepackage[american, nooldvoltagedirection, siunitx]{circuitikz}

% Unidades (inclusa no circuitikz)
%\usepackage{siunitx}

\usepackage{tikz}
\usetikzlibrary{positioning, calc,
patterns, arrows.meta, decorations.pathmorphing,
circuits.plc.ladder}

\newlength{\ladderskip}
\setlength{\ladderskip}{5\tikzcircuitssizeunit} % 5\tikzcircuitssizeunit = 35pt
\newlength{\ladderrungsep}
\setlength{\ladderrungsep}{.2\ladderskip}
\def\ladderrungend#1{\pgftransformyshift{-#1\ladderskip-\ladderrungsep}}
\newcommand{\powerrails}[1][0.7]{\draw let \p1=(laddertopright) in
	(0,\y1+#1\ladderskip) -- (0,\ladderskip)
	(\x1,\y1+#1\ladderskip) -- (\x1,\ladderskip);}

% Cores (inclusa no tikz)
%\usepackage{xcolor}


%%%%%%%%%%%%%%%%%%%%%%%%%%%%%%%%%%%%%%%%%%%
%%%%%%%%%%%%%%%%%%%%%%%%%%%%%%%%%%%%%%%%%%%
%%%%%%%%%%%%%%%%% COISAS DOIDAS %%%%%%%%%%%%%%%%%

% setup do siunitx
\sisetup{per-mode=symbol,output-decimal-marker={,},binary-units}

% MATLAB
%\usepackage{xspace}
%\newcommand{\MATLAB}{\textsc{Matlab}\xspace}

% Notação de negação lógica
\newcommand{\notted}[1]{%
	\overline{#1}%
}

% Padrões de linhas que não funcionam no overleaf

%\pgfdeclarepatternformonly{south east lines}{\pgfqpoint{-0pt}{-0pt}}{\pgfqpoint{3pt}{3pt}}{\pgfqpoint{3pt}{3pt}}{
%	\pgfsetlinewidth{0.4pt}
%	\pgfpathmoveto{\pgfqpoint{0pt}{3pt}}
%	\pgfpathlineto{\pgfqpoint{3pt}{0pt}}
%	\pgfpathmoveto{\pgfqpoint{.2pt}{-.2pt}}
%	\pgfpathlineto{\pgfqpoint{-.2pt}{.2pt}}
%	\pgfpathmoveto{\pgfqpoint{3.2pt}{2.8pt}}
%	\pgfpathlineto{\pgfqpoint{2.8pt}{3.2pt}}
%	\pgfusepath{stroke}}
%
%\pgfdeclarepatternformonly{south west lines}{\pgfqpoint{-0pt}{-0pt}}{\pgfqpoint{3pt}{3pt}}{\pgfqpoint{3pt}{3pt}}{
%	\pgfsetlinewidth{0.4pt}
%	\pgfpathmoveto{\pgfqpoint{0pt}{0pt}}
%	\pgfpathlineto{\pgfqpoint{3pt}{3pt}}
%	\pgfpathmoveto{\pgfqpoint{2.8pt}{-.2pt}}
%	\pgfpathlineto{\pgfqpoint{3.2pt}{.2pt}}
%	\pgfpathmoveto{\pgfqpoint{-.2pt}{2.8pt}}
%	\pgfpathlineto{\pgfqpoint{.2pt}{3.2pt}}
%	\pgfusepath{stroke}}


% cor do fundo do block
%\definecolor{mWhite}{RGB}{239, 230, 231}

% Definições de funções trinométricas em pt-br
%\DeclareMathOperator{\sen}{sen}
%\DeclareMathOperator{\tg}{tg}
%\DeclareMathOperator{\cotg}{cotg}
%\DeclareMathOperator{\cossec}{cossec}

% QED branco
%\newcommand*{\QEDB}{\hfill\ensuremath{\square}}

% estilos para diagrama de blocos
%\newcommand{\deftkzbds}{
%	\tikzstyle{block} = [draw, fill=blue!20, rectangle, minimum height=3em, minimum width=6em]
%	\tikzstyle{sum} = [draw, fill=blue!20, circle, node distance=1cm]
%	\tikzstyle{input} = [coordinate]
%	\tikzstyle{output} = [coordinate]
%	\tikzstyle{pinstyle} = [pin edge={to-,thin,black}]
%}

% tipo de coluna matemática centralizada
%\newcolumntype{C}{>{$}c<{$}}

% para marcar ponto na tela e desenhar sobre
\newcommand{\tikzmark}[1]{\tikz[baseline,remember picture] \coordinate (#1) {};}

% unidades uteis para siuntix

\DeclareSIUnit{\lbf}{lbf}
\DeclareSIUnit{\kgf}{kgf}
\DeclareSIUnit{\kgfp}{\kgf \per \centi\meter\squared}
\DeclareSIUnit{\mca}{mca}
\DeclareSIUnit{\barp}{bar}
\DeclareSIUnit{\pol}{pol}
\DeclareSIUnit{\psip}{psi}
\DeclareSIUnit{\atm}{atm}
\DeclareSIUnit{\HP}{HP}
\DeclareSIUnit{\psid}{\lbf\per\pol\squared}
%\DeclareSIUnit{}{}

% Usar para diminuir espaçamento antes/depois da equação (\useshortskip)
%\usepackage{nccmath}
%\usepackage{xpatch}
%\xpatchcmd{\NCC@ignorepar}{%
%	\abovedisplayskip\abovedisplayshortskip}
%{%
%	\abovedisplayskip\abovedisplayshortskip%
%	\belowdisplayskip\belowdisplayshortskip}
%{}{}

\usepackage{chngcntr}
\counterwithin*{equation}{section}
\newcounter{saveenumi}
\newcommand{\saveenumerate}{%
	\stepcounter{saveenumi}%
	\label{saveenumi-\thesaveenumi}}
\newcommand{\restoreenumerate}{%
	\setcounterref{enumi}{saveenumi-\thesaveenumi}}

% Símbolo para graus (usar \ang{degrees} do package siunitx)
%\newcommand{\degree}{\ensuremath{^\circ}}

% X de errado e certinho correspondente
\usepackage{pifont}
\newcommand{\cmark}{\ding{51}}%
\newcommand{\xmark}{\ding{55}}%

% \itemequation[label]{text before}{equation}

%\makeatletter
%\newcommand*{\itemequation}[3][]{%
%  \item
%  \begingroup
%    \refstepcounter{equation}%
%    \ifx\\#1\\%
%    \else
%      \label{#1}%
%    \fi
%    \sbox0{#2}%
%    \sbox2{$\displaystyle#3\m@th$}%
%    \sbox4{ \@eqnnum}%
%    \dimen@=.5\dimexpr\linewidth-\wd2\relax
%    % Warning for overlapping
%    \let\CenterInSpace=N%
%    \ifcase
%    \ifdim\wd0>\dimen@
%          \z@
%        \else
%          \ifdim\wd4>\dimen@
%            \z@
%          \else
%            \@ne
%          \fi
%        \fi
%      \let\CenterInSpace=Y%
%    \fi
%    \ifdim\dimexpr\wd0+\wd2+\wd4\relax>\linewidth
%      \@latex@warning{Equation is too large}%
%    \fi
%    \noindent
%    \rlap{\copy0}%
%    \ifx\CenterInSpace Y%
%      \rlap{\hbox to \linewidth{\kern\wd0\hss\copy2\hss\kern\wd4}}%
%    \else
%      \rlap{\hbox to \linewidth{\hfill\copy2\hfill}}%
%    \fi
%    \hbox to \linewidth{\hfill\copy4}%
%    \hspace{0pt}% allow linebreak
%  \endgroup
%  \ignorespaces
%}
%\makeatother

%%%%%%%%%%%%%%%%%%%%%%%%%%%%%%%%%%%%%%%%%%%
%%%%%%%%%%%%%%%%%%%%%%%%%%%%%%%%%%%%%%%%%%%
%%%%%%%%%%%%%%%%%% RODAPÉ %%%%%%%%%%%%%%%%%

\setbeamercolor{footline}{fg=white}
\setbeamertemplate{footline}
{\begin{tikzpicture}
    \node [inner sep=0pt, anchor=east] (0,0) {\includegraphics[width=\paperwidth,height=1cm]{Figuras/Capa/macaefooter.png}};
    \node [inner sep=0pt, anchor=east] at (-2ex,-3ex) {\insertframenumber{} / \inserttotalframenumber};
\end{tikzpicture}}


%%%%%%%%%%%%%%%%%%%%%%%%%%%%%%%%%%%%%%%%%%%
%%%%%%%%%%%%%%%%%%%%%%%%%%%%%%%%%%%%%%%%%%%
%%%%%%%%%%% INFORMAÇÕES DO CURSO %%%%%%%%%%

\title[Princípios de Automação]
{Princípios de Automação}

\subtitle {Notas de aula}

\author{Prof. Yago Pessanha Corrêa}

\institute[MSP/IFF] 
{
Laboratório de Mecatrônica e Processamento de Sinais (MSP) \\
Instituto Federal de Educação, Ciência e Tecnologia Fluminense (IFFluminense) \\
Curso Técnico Subsequente em Eletrônica \\
\vspace*{.1cm} {\tt \textbf{yago.correa@iff.edu.br}}\\
}

%\tddate


%%%%%%%%%%%%%%%%%%%%%%%%%%%%%%%%%%%%%%%%%%%
%%%%%%%%%%%%%%%%%%%%%%%%%%%%%%%%%%%%%%%%%%%
%%%%%%%%%%%%%%%%% SUMÁRIO %%%%%%%%%%%%%%%%%

\AtBeginSection[]
{
  \begin{frame}<beamer>{Sumário}
    \tableofcontents[currentsection]
  \end{frame}
}